\documentclass{dlproblemset}
\usepackage{tikz}
\usepackage{emoji}
\tikzset{every picture/.style={line width=0.75pt}}

\title{Tokenização}
\subtitle{Lista de Exercícios}

\begin{document}

\begin{enumerate}[I)]

% ------------------------------------------------------------------
% Q1 — Fácil
% ------------------------------------------------------------------
\item Quais são os três componentes fundamentais de um sistema de tokenização?

\begin{enumerate}[a)]
    \item Pré-processamento, normalização e lematização.
    \item Vocabulário, segmentação e mapeamento.
    \item Embedding, atenção e decodificação.
    \item Codificação, compressão e indexação.
\end{enumerate}

% ------------------------------------------------------------------
% Q2 — Médio
% ------------------------------------------------------------------
\item Em tokenizers com pré-tokenização, como os do GPT-2 e do BERT, nenhum token
      do vocabulário contém duas palavras (como \texttt{of the}). O que garante
      essa propriedade?

\begin{enumerate}[a)]
    \item O espaço possui um token próprio que nunca participa de merges.
    \item Tokens com espaço interno são descartados na filtragem final do vocabulário.
    \item A contagem e a aplicação dos merges ocorrem apenas dentro das fronteiras
          dos pré-tokens.
    \item O marcador de fim de palavra impede que o último símbolo de uma palavra
          forme pares.
\end{enumerate}

% ------------------------------------------------------------------
% Q3 — Médio
% ------------------------------------------------------------------
\item No tokenizer do GPT-2, o espaço não é descartado: ele é embutido no token
      seguinte (ex.: \texttt{\textvisiblespace world}). Qual é a consequência dessa
      convenção para uma mesma palavra que aparece ora no início, ora no meio de
      uma frase?

\begin{enumerate}[a)]
    \item O tokenizer normaliza os dois casos para o mesmo token, e o espaço é
          restaurado por regras de detokenização.
    \item Os dois casos compartilham o mesmo ID, mas recebem positional encodings
          diferentes.
    \item O vocabulário dobra de tamanho, pois toda palavra precisa registrar as
          duas variantes.
    \item Os dois casos correspondem a tokens com IDs e embeddings distintos.
\end{enumerate}

% ------------------------------------------------------------------
% Q4 — Fácil
% ------------------------------------------------------------------
\item O algoritmo BPE (Sennrich et al., 2016) constrói o vocabulário em qual direção
      e utiliza qual critério para selecionar os pares a mesclar?

\begin{enumerate}[a)]
    \item De cima para baixo (top-down), removendo tokens com menor probabilidade.
    \item De baixo para cima (bottom-up), mesclando iterativamente o par de símbolos
          adjacentes mais frequente no corpus.
    \item De baixo para cima (bottom-up), mesclando o par que maximiza a verossimilhança
          do modelo de linguagem unigrama.
    \item De cima para baixo (top-down), particionando tokens longos nos pontos de
          menor frequência de bigrama.
\end{enumerate}

% ------------------------------------------------------------------
% Q5 — Difícil (inferência do BPE)
% ------------------------------------------------------------------
\item A tabela de merges do BPE aprendida na aula, a partir do corpus
      \texttt{low} (×5), \texttt{lower} (×2), \texttt{newest} (×6), \texttt{widest} (×3),
      é, em ordem:
      \begin{center}
      1.\ \texttt{(e,s)} $\rightarrow$ \texttt{es} \quad
      2.\ \texttt{(es,t)} $\rightarrow$ \texttt{est} \quad
      3.\ \texttt{(est,\_)} $\rightarrow$ \texttt{est\_} \quad
      4.\ \texttt{(l,o)} $\rightarrow$ \texttt{lo} \\
      5.\ \texttt{(lo,w)} $\rightarrow$ \texttt{low} \quad
      6.\ \texttt{(n,e)} $\rightarrow$ \texttt{ne} \quad
      7.\ \texttt{(ne,w)} $\rightarrow$ \texttt{new}
      \end{center}
      Aplicando esses merges na inferência, como é codificada a palavra \emph{nova}
      \texttt{slowest} (divisão inicial: \texttt{s~l~o~w~e~s~t~\_})?

\begin{enumerate}[a)]
    \item \texttt{slow}~\texttt{est\_}
    \item \texttt{s}~\texttt{low}~\texttt{est\_}
    \item \texttt{s}~\texttt{lo}~\texttt{w}~\texttt{est\_}
    \item \texttt{s}~\texttt{low}~\texttt{est}~\texttt{\_}
\end{enumerate}

% ------------------------------------------------------------------
% Q6 — Médio
% ------------------------------------------------------------------
\item No \textbf{Byte-level BPE} (adotado pelo GPT-2), qual é o tamanho do vocabulário
      base \emph{antes} de qualquer merge ser aplicado, e por quê esse tamanho garante
      que nenhum token seja desconhecido (\texttt{<unk>})?

\begin{enumerate}[a)]
    \item 65.536 tokens (um por ponto de código Unicode BMP), cobrindo os caracteres
          mais comuns de todas as línguas.
    \item 256 tokens (um por valor de byte possível), pois qualquer texto pode ser
          codificado em UTF-8 como uma sequência de bytes no intervalo $[0, 255]$.
    \item 128 tokens (caracteres ASCII), suficientes para texto em inglês sem OOV.
    \item 1.024 tokens (pares de bytes comuns), reduzindo o comprimento das sequências
          sem perder cobertura.
\end{enumerate}

% ------------------------------------------------------------------
% Q7 — Difícil (cálculo de score WordPiece)
% ------------------------------------------------------------------
\item Considere o corpus de treinamento \texttt{low} (×5), \texttt{lower} (×2),
      \texttt{newest} (×6), \texttt{widest} (×3). O WordPiece seleciona merges
      maximizando o score:
      \[
        \mathrm{score}(x, y) = \frac{\mathrm{freq}(xy)}{\mathrm{freq}(x) \cdot \mathrm{freq}(y)}.
      \]
      Calcule os scores para os pares \texttt{(e,s)} e \texttt{(i,d)}.
      Qual par o WordPiece mescla \emph{primeiro}? Esse resultado é igual ao BPE?

\begin{enumerate}[a)]
    \item \texttt{(e,s)}: score $\approx 0.059$; \texttt{(i,d)}: score $\approx 0.333$.
          O WordPiece mescla \texttt{(i,d)} primeiro; o BPE mesclaria \texttt{(e,s)}.
    \item \texttt{(e,s)}: score $\approx 0.529$; \texttt{(i,d)}: score $\approx 0.111$.
          O WordPiece mescla \texttt{(e,s)} primeiro; o BPE também mesclaria \texttt{(e,s)}.
    \item \texttt{(e,s)}: score $= 9$; \texttt{(i,d)}: score $= 3$.
          O WordPiece mescla \texttt{(e,s)} primeiro; o BPE também mesclaria \texttt{(e,s)}.
    \item \texttt{(e,s)}: score $\approx 0.059$; \texttt{(i,d)}: score $\approx 0.333$.
          O WordPiece mescla \texttt{(e,s)} primeiro; o BPE mesclaria \texttt{(i,d)}.
\end{enumerate}

% ------------------------------------------------------------------
% Q8 — Médio (inferência do WordPiece)
% ------------------------------------------------------------------
\item Na inferência, o WordPiece usa longest-match guloso da esquerda para a direita.
      Com um vocabulário que contém, entre outros, os tokens \texttt{un},
      \texttt{happy}, \texttt{\#\#happi}, \texttt{\#\#ness} e \texttt{\#\#ly},
      qual é a codificação da palavra \texttt{unhappiness}?

\begin{enumerate}[a)]
    \item \texttt{un}~\texttt{happy}~\texttt{\#\#ness}
    \item \texttt{unhappi}~\texttt{\#\#ness}
    \item \texttt{un}~\texttt{\#\#happi}~\texttt{\#\#ness}
    \item \texttt{un}~\texttt{\#\#happy}~\texttt{\#\#ness}
\end{enumerate}

% ------------------------------------------------------------------
% Q9 — Médio
% ------------------------------------------------------------------
\item Durante a inferência do WordPiece, nenhuma combinação de tokens do vocabulário
      consegue cobrir uma palavra da entrada. O que acontece com essa palavra, e o
      que aconteceria com o byte-level BPE na mesma situação?

\begin{enumerate}[a)]
    \item No WordPiece, só o trecho não coberto vira \texttt{[UNK]}; no byte-level
          BPE, o mesmo trecho vira \texttt{<unk>}.
    \item No WordPiece, a palavra é decomposta em caracteres; no byte-level BPE,
          em pontos de código Unicode.
    \item No WordPiece, a palavra é truncada no maior prefixo coberto; no byte-level
          BPE, o restante é completado com bytes.
    \item No WordPiece, a palavra inteira vira \texttt{[UNK]}; no byte-level BPE,
          ela é decomposta até o nível de bytes.
\end{enumerate}

% ------------------------------------------------------------------
% Q10 — Médio
% ------------------------------------------------------------------
\item O tokenizador Unigram Language Model (Kudo, 2018) difere do BPE principalmente
      em qual aspecto do processo de construção do vocabulário?

\begin{enumerate}[a)]
    \item O Unigram usa merges guiados por frequência de bigramas, enquanto o BPE
          usa verossimilhança.
    \item O Unigram é top-down: parte de um vocabulário grande e itera removendo os
          tokens que menos contribuem para a verossimilhança do corpus.
    \item O Unigram aplica dropout nos merges durante o treinamento, enquanto o BPE
          usa merges determinísticos.
    \item O Unigram é restrito a modelos que usam SentencePiece, enquanto o BPE pode
          ser usado com qualquer framework.
\end{enumerate}

% ------------------------------------------------------------------
% Q11 — Médio
% ------------------------------------------------------------------
\item O \textbf{BPE Dropout} (Provilkov et al., 2020) e a \textbf{subword
      regularization} do Unigram (Kudo, 2018) compartilham qual objetivo de
      treinamento?

\begin{enumerate}[a)]
    \item Reduzir o tamanho do vocabulário durante o treinamento para diminuir o
          uso de memória.
    \item Aumentar a velocidade de tokenização ao pular merges pouco frequentes.
    \item Servir como \emph{data augmentation}, expondo o modelo a múltiplas
          segmentações válidas de um mesmo texto durante o treinamento.
    \item Garantir que cada token apareça no mínimo uma vez por batch de treinamento.
\end{enumerate}

% ------------------------------------------------------------------
% Q12 — Fácil
% ------------------------------------------------------------------
\item Por que o SentencePiece (Kudo \& Richardson, 2018) é descrito como
      \emph{independente de idioma} (language-agnostic)?

\begin{enumerate}[a)]
    \item Porque suporta apenas o algoritmo BPE, aplicável de forma idêntica a
          qualquer corpus de treinamento.
    \item Porque utiliza um vocabulário base fixo de 256 bytes, idêntico para
          todos os sistemas de escrita.
    \item Porque dispensa a pré-tokenização por espaços em branco, tratando o
          espaço como um caractere regular do fluxo de entrada.
    \item Porque aprende um vocabulário separado para cada idioma detectado
          automaticamente no corpus.
\end{enumerate}

% ------------------------------------------------------------------
% Q13 — Médio
% ------------------------------------------------------------------
\item Qual alternativa descreve corretamente o papel das bibliotecas SentencePiece
      e tiktoken no ecossistema de tokenização?

\begin{enumerate}[a)]
    \item O SentencePiece apenas aplica vocabulários prontos do Google; o tiktoken
          treina novos vocabulários a partir de qualquer corpus.
    \item Ambos treinam vocabulários; o SentencePiece é restrito ao Unigram e o
          tiktoken, ao WordPiece.
    \item Ambos apenas aplicam vocabulários prontos; a diferença está no formato
          de serialização da tabela de merges.
    \item O SentencePiece treina e aplica vocabulários com BPE ou Unigram; o
          tiktoken apenas aplica vocabulários byte-level BPE prontos.
\end{enumerate}

% ------------------------------------------------------------------
% Q14 — Médio
% ------------------------------------------------------------------
\item Sobre tokens especiais como \texttt{[CLS]}, \texttt{<|endoftext|>} e
      \texttt{<|im\_start|>}: qual propriedade os define, e qual risco surge quando
      texto do usuário que os imita não é escapado?

\begin{enumerate}[a)]
    \item São atômicos, nunca resultando de merges nem sendo divididos; sem escape,
          documentos podem injetar marcadores estruturais na conversa.
    \item São aprendidos como os demais tokens, com IDs reservados no fim do
          vocabulário; sem escape, o texto é truncado no primeiro marcador.
    \item São expandidos em sequências de bytes na entrada; sem escape, a
          detokenização perde os caracteres \texttt{<} e \texttt{>}.
    \item São ignorados pela camada de embedding; sem escape, o modelo processa o
          texto normalmente, sem risco prático.
\end{enumerate}

% ------------------------------------------------------------------
% Q15 — Médio
% ------------------------------------------------------------------
\item Tokens como \texttt{\textvisiblespace SolidGoldMagikarp} faziam o GPT-3
      responder de forma errática a pedidos simples como ``repita esta palavra''.
      Qual é a origem desse fenômeno (\emph{glitch tokens})?

\begin{enumerate}[a)]
    \item O token colidia com um token especial de controle, e o modelo o
          interpretava como marcador de fim de documento.
    \item O token continha bytes fora do intervalo UTF-8 válido, corrompendo o
          mapeamento byte-para-unicode do tokenizer.
    \item O vocabulário foi aprendido num corpus diferente do usado para treinar o
          modelo, e o embedding do token quase não recebeu gradiente.
    \item O token era longo demais e foi truncado em bytes inválidos durante a
          codificação da entrada.
\end{enumerate}

% ------------------------------------------------------------------
% Q16 — Difícil (custo de embeddings)
% ------------------------------------------------------------------
\item Um modelo de linguagem tem vocabulário de tamanho $|V| = 50.000$ e dimensão
      de embedding $d = 1.024$, \emph{sem} weight tying entre a tabela de embedding
      e a camada de saída (\texttt{lm\_head}). Quantos parâmetros essas duas
      matrizes somam, e o que muda com weight tying?

\begin{enumerate}[a)]
    \item Somam 102,4M de parâmetros; com weight tying, a mesma matriz é reutilizada
          na saída e o total cai para 51,2M.
    \item Somam 51,2M de parâmetros; com weight tying, o total cai para 25,6M.
    \item Somam 102,4M de parâmetros; com weight tying, o total não muda, pois
          apenas os gradientes são compartilhados.
    \item Somam 51,2M de parâmetros; com weight tying, o total dobra para 102,4M,
          pois a matriz é duplicada na saída.
\end{enumerate}

\end{enumerate}


% ============================================================
\section*{Soluções}
% ============================================================

% --- Solução Q1 ---
\subsection*{I — Resposta: (b)}
Um sistema de tokenização é definido por três componentes:
\begin{enumerate}
  \item \textbf{Vocabulário}: o conjunto de tokens reconhecidos pelo modelo, ou seja,
        quais sequências de caracteres são unidades atômicas.
  \item \textbf{Segmentação}: o algoritmo que divide um texto de entrada em uma sequência
        de tokens do vocabulário.
  \item \textbf{Mapeamento}: a função que converte cada token em um ID inteiro (e vice-versa),
        necessária para alimentar o modelo com tensores numéricos.
\end{enumerate}
Essas três decisões em conjunto determinam silenciosamente o desempenho do modelo,
a eficiência computacional e a equidade multilíngue.
% ---

% --- Solução Q2 ---
\subsection*{II — Resposta: (c)}
A pré-tokenização é o ``passo zero'' do pipeline: regras fixas (espaço, pontuação
ou uma regex) dividem o texto em pré-tokens \emph{antes} de qualquer algoritmo de
subword. Tanto o treinamento (contagem de pares) quanto a inferência (aplicação dos
merges) operam \textbf{dentro} dessas fronteiras, então nenhum merge pode unir
símbolos de palavras diferentes e nenhum token como \texttt{of the} pode surgir.
Uma consequência prática importante: o corpus de treinamento pode ser reduzido a
uma tabela de palavras com frequências, exatamente o formato do exemplo recorrente
da aula. O SentencePiece é a exceção: dispensa essa etapa e trata o espaço como um
caractere comum do fluxo de entrada.
% ---

% --- Solução Q3 ---
\subsection*{III — Resposta: (d)}
Embutir o espaço no token seguinte torna a tokenização \textbf{reversível} (o
espaço não é descartado), mas cria duas entradas de vocabulário para a mesma
palavra: \texttt{world} (início de frase ou após pontuação) e
\texttt{\textvisiblespace world} (meio de frase). São tokens com IDs diferentes e,
portanto, com \textbf{embeddings diferentes}: o modelo precisa aprender por conta
própria que os dois se referem à mesma palavra. Internamente o GPT-2 grafa o espaço
embutido como \texttt{\.{G}} (ex.: \texttt{\.{G}world}), um artefato do mapeamento
byte-para-unicode. As demais alternativas descrevem mecanismos que não existem
nessa convenção: não há normalização dos dois casos, o ID é distinto (não apenas a
posição) e o vocabulário não registra ambas as variantes de toda palavra, apenas
das que aparecem nos dois contextos com frequência suficiente.
% ---

% --- Solução Q4 ---
\subsection*{IV — Resposta: (b)}
O BPE (Byte Pair Encoding) é um algoritmo \textbf{bottom-up} e \textbf{guloso}:
\begin{enumerate}
  \item Inicia com um vocabulário de símbolos individuais (caracteres ou bytes).
  \item Conta a frequência de todos os pares de símbolos adjacentes no corpus.
  \item Mescla o par mais frequente em um novo símbolo composto.
  \item Repete o processo por $k$ iterações (hiperparâmetro que define o tamanho do vocabulário).
\end{enumerate}
A tabela de merges aprendida é aplicada deterministicamente na inferência, na mesma
ordem em que foi aprendida. O BPE foi proposto para tradução automática neural por
Sennrich et al.\ (2016) e tornou-se a base dos tokenizadores da família GPT.
% ---

% --- Solução Q5 ---
\subsection*{V — Resposta: (b)}
Na inferência, os merges são aplicados \textbf{na ordem em que foram aprendidos},
sempre que o par correspondente estiver presente:
\begin{enumerate}
  \item Início: \texttt{s l o w e s t \_}
  \item Merge 1, \texttt{(e,s)} $\rightarrow$ \texttt{es}: \quad \texttt{s l o w es t \_}
  \item Merge 2, \texttt{(es,t)} $\rightarrow$ \texttt{est}: \quad \texttt{s l o w est \_}
  \item Merge 3, \texttt{(est,\_)} $\rightarrow$ \texttt{est\_}: \quad \texttt{s l o w est\_}
  \item Merge 4, \texttt{(l,o)} $\rightarrow$ \texttt{lo}: \quad \texttt{s lo w est\_}
  \item Merge 5, \texttt{(lo,w)} $\rightarrow$ \texttt{low}: \quad \texttt{s low est\_}
  \item Merges 6 e 7 não se aplicam (não há \texttt{(n,e)} nem \texttt{(ne,w)}).
\end{enumerate}
Resultado: \texttt{s}~\texttt{low}~\texttt{est\_}, com 3 tokens. Note que
\texttt{slow} não pode surgir: o par \texttt{(s,low)} nunca foi aprendido como
merge. Mesmo sem ``slowest'' jamais ter aparecido no corpus, a palavra é composta
por subwords aprendidos, compartilhados com ``low'', ``newest'' e ``widest''.
% ---

% --- Solução Q6 ---
\subsection*{VI — Resposta: (b)}
O UTF-8 codifica qualquer ponto de código Unicode em uma sequência de 1 a 4 bytes,
onde cada byte assume um valor no intervalo $[0, 255]$. Ao operar sobre bytes brutos,
o byte-level BPE começa com exatamente \textbf{256 tokens base}, um para cada valor
de byte possível. Consequência direta: qualquer texto, em qualquer idioma ou sistema
de escrita, incluindo emojis e caracteres especiais, pode ser representado sem
nenhum token desconhecido (\texttt{<unk>}). Os merges do BPE são então aplicados sobre
esses 256 bytes para construir tokens maiores. O GPT-2 usa esse esquema com um
vocabulário final de 50.257 tokens (256 bytes + 50.000 merges + 1 token especial (\texttt{<EOS>})).
O GPT-3 reutiliza essencialmente o mesmo tokenizer (\texttt{p50k\_base}, 50.281 tokens), enquanto
o GPT-4 adota um vocabulário maior, o \texttt{cl100k\_base}, com 100.256 tokens, melhorando
a compressão (menos tokens por texto) e a cobertura multilíngue.
% ---

% --- Solução Q7 ---
\subsection*{VII — Resposta: (a)}

\textbf{Frequências individuais dos caracteres} (calculadas somando sobre todas as
ocorrências no corpus):
\begin{itemize}
  \item \texttt{freq(e)} $= 1{\times}2 + 2{\times}6 + 1{\times}3 = 2 + 12 + 3 = 17$
        \quad (``lower'': 1 \texttt{e}; ``newest'': 2 \texttt{e}'s; ``widest'': 1 \texttt{e})
  \item \texttt{freq(s)} $= 6 + 3 = 9$
  \item \texttt{freq(i)} $= 3$ \quad (só em ``widest'')
  \item \texttt{freq(d)} $= 3$ \quad (só em ``widest'')
\end{itemize}

\textbf{Frequências dos pares:}
\begin{itemize}
  \item \texttt{freq(es)} $= 6 + 3 = 9$
  \item \texttt{freq(id)} $= 3$
\end{itemize}

\textbf{Scores WordPiece:}
\[
  \mathrm{score}(e, s)
  = \frac{\mathrm{freq}(es)}{\mathrm{freq}(e) \cdot \mathrm{freq}(s)}
  = \frac{9}{17 \times 9}
  = \frac{9}{153}
  \approx 0.059.
\]
\[
  \mathrm{score}(i, d)
  = \frac{\mathrm{freq}(id)}{\mathrm{freq}(i) \cdot \mathrm{freq}(d)}
  = \frac{3}{3 \times 3}
  = \frac{3}{9}
  \approx 0.333.
\]

\textbf{Conclusão:} o WordPiece mescla \texttt{(i,d)} primeiro (score $\approx 0.333$),
embora \texttt{(e,s)} tenha frequência bruta três vezes maior (9 vs 3).

\textbf{Intuição:} \texttt{i} e \texttt{d} coocorrem quase exclusivamente juntos no
corpus: toda vez que \texttt{i} aparece, \texttt{d} é o próximo símbolo. O score
captura essa alta informação mútua. Por outro lado, \texttt{e} e \texttt{s} aparecem
frequentemente de forma \emph{independente} em outros contextos, reduzindo o score
mesmo com alta frequência de coocorrência absoluta. O BPE ignora essa distinção e
escolheria \texttt{(e,s)} simplesmente por ter maior contagem.
% ---

% --- Solução Q8 ---
\subsection*{VIII — Resposta: (c)}
O longest-match guloso varre a palavra da esquerda para a direita, sempre tomando o
maior token do vocabulário que casa com o início do trecho restante:
\begin{enumerate}
  \item Trecho \texttt{unhappiness}: a maior correspondência inicial é \texttt{un}.
  \item Trecho restante \texttt{happiness}: como não é início de palavra, apenas
        tokens de continuação (prefixo \texttt{\#\#}) são válidos. A maior
        correspondência é \texttt{\#\#happi}. Note que \texttt{happy} não serve por
        duas razões: não tem o prefixo \texttt{\#\#} e não é prefixo do trecho
        (\texttt{happy} $\neq$ \texttt{happi}$\ldots$).
  \item Trecho restante \texttt{ness}: maior correspondência \texttt{\#\#ness}.
\end{enumerate}
Resultado: \texttt{un}~\texttt{\#\#happi}~\texttt{\#\#ness}. O prefixo \texttt{\#\#}
indica que o token continua a palavra anterior, permitindo distinguir, por exemplo,
o \texttt{ness} que inicia uma palavra do \texttt{\#\#ness} que a encerra.
% ---

% --- Solução Q9 ---
\subsection*{IX — Resposta: (d)}
No WordPiece, se nenhuma combinação de tokens do vocabulário cobre a palavra, ela
\textbf{inteira} é substituída pelo token especial \texttt{[UNK]}, e todo o seu
conteúdo se perde: o modelo recebe apenas ``havia aqui uma palavra desconhecida''.
O byte-level BPE, em contraste, nunca produz tokens desconhecidos: qualquer texto é
representável como sequência de bytes UTF-8, e o pior caso é uma decomposição em
tokens de 1 byte (sequência mais longa, mas sem perda de informação). O BPE padrão
em nível de caractere e o Unigram também degradam graciosamente, decompondo até os
símbolos individuais do vocabulário base, razão pela qual caracteres individuais
nunca são podados no Unigram.
% ---

% --- Solução Q10 ---
\subsection*{X — Resposta: (b)}
Enquanto o BPE é \textbf{bottom-up} (começa pequeno e cresce mesclando), o Unigram
é \textbf{top-down} (começa grande e encolhe podando):
\begin{enumerate}
  \item Inicializa o vocabulário com todos os substrings do corpus até um comprimento
        máximo (podendo ter centenas de milhares de entradas).
  \item Define uma distribuição de probabilidade unigrama sobre o vocabulário.
  \item Para cada token candidato à remoção, calcula a queda na verossimilhança do
        corpus se ele fosse eliminado.
  \item Remove os tokens com menor queda (os menos úteis), tipicamente 10–20\% do
        vocabulário por iteração.
  \item Repete até atingir o tamanho de vocabulário desejado.
\end{enumerate}
Uma propriedade exclusiva do Unigram é poder gerar \emph{múltiplas} segmentações
válidas de um mesmo texto com probabilidades associadas, algo explorado pela
subword regularization como augmentation de dados.
% ---

% --- Solução Q11 ---
\subsection*{XI — Resposta: (c)}
Ambas as técnicas introduzem \textbf{aleatoriedade controlada} na segmentação durante
o treinamento:
\begin{itemize}
  \item \textbf{BPE Dropout} (Provilkov et al., 2020): durante o treinamento, cada
        merge do BPE é ignorado com probabilidade $p$ (tipicamente $p = 0.1$).
        A mesma palavra pode ser tokenizada de maneiras diferentes em diferentes passos,
        como \texttt{low~est\_}, \texttt{lo~w~est\_} ou \texttt{l~o~w~est\_}.
  \item \textbf{Subword regularization} (Kudo, 2018): usa o modelo probabilístico do
        Unigram para amostrar uma segmentação dentre as válidas, em vez de sempre
        escolher a de maior probabilidade.
\end{itemize}
O efeito é equivalente ao dropout nas camadas de rede: força o modelo a ser robusto
a múltiplas representações de uma mesma sequência, melhorando a generalização,
especialmente em cenários com dados escassos ou idiomas de alta morfologia.
% ---

% --- Solução Q12 ---
\subsection*{XII — Resposta: (c)}
A maioria dos tokenizadores depende de regras de pré-tokenização baseadas em espaços
em branco e pontuação, regras que funcionam bem para o inglês mas são inadequadas para
idiomas como japonês, chinês ou tailandês, que não usam espaços para delimitar palavras.
O SentencePiece elimina essa dependência ao tratar o texto de entrada como um fluxo
bruto de bytes ou pontos de código Unicode, sem nenhuma segmentação prévia. O espaço
em branco é substituído pelo símbolo \texttt{\_} e tratado como qualquer outro caractere,
tornando a tokenização completamente reversível e uniforme em todos os idiomas.
O SentencePiece serve como \emph{framework}: internamente pode usar BPE ou o modelo
Unigram como algoritmo de segmentação.
% ---

% --- Solução Q13 ---
\subsection*{XIII — Resposta: (d)}
As duas bibliotecas têm papéis diferentes no ecossistema:
\begin{itemize}
  \item \textbf{SentencePiece} (Google, 2018) cobre \textbf{treino e inferência}: o
        usuário treina seu próprio vocabulário sobre um corpus, escolhendo BPE ou
        Unigram como algoritmo, com normalização NFKC configurável e
        \texttt{byte\_fallback} opcional para evitar \texttt{<unk>}. É usado por
        T5, LLaMA 1/2, Mistral e Gemma.
  \item \textbf{tiktoken} (OpenAI, 2022) é uma biblioteca de \textbf{inferência
        apenas}: aplica vocabulários byte-level BPE já prontos
        (\texttt{cl100k\_base}, \texttt{o200k\_base}, \ldots) com pré-tokenização
        obrigatória via regex. É usado pelos modelos GPT-3.5/4/4o e pelo LLaMA 3.
\end{itemize}
O HuggingFace \texttt{tokenizers} reimplementa ambos os formatos e é o que o
\texttt{AutoTokenizer} carrega por baixo dos panos.
% ---

% --- Solução Q14 ---
\subsection*{XIV — Resposta: (a)}
Tokens especiais são \textbf{atômicos}: são inseridos diretamente no vocabulário,
nunca resultam de merges e nunca são divididos pela pré-tokenização. Eles carregam
\emph{estrutura}, não conteúdo: \texttt{[CLS]} agrega a representação para
classificação no BERT, \texttt{<|endoftext|>} marca fim de documento no GPT-2, e
pares como \texttt{<|im\_start|>}/\texttt{<|im\_end|>} delimitam turnos e papéis
(user, assistant) nos chat templates. Justamente por carregarem estrutura, texto do
usuário que imita um token especial precisa ser \textbf{escapado}: se a string
\texttt{<|im\_start|>} de um documento fosse codificada como o token especial, o
documento injetaria um marcador de turno na conversa (\emph{prompt injection}).
No tiktoken, por padrão, tokens especiais encontrados no texto do usuário geram
\textbf{erro} (\texttt{disallowed\_special}), deixando o escape como
responsabilidade explícita da aplicação.
% ---

% --- Solução Q15 ---
\subsection*{XV — Resposta: (c)}
O vocabulário BPE do GPT-2/3 foi aprendido num corpus \textbf{diferente} do usado
para treinar o modelo. Strings como \texttt{\textvisiblespace SolidGoldMagikarp}
(usernames do Reddit) eram frequentes o bastante no corpus do tokenizer para
ganharem um token próprio, mas foram filtradas do corpus de treinamento do modelo.
Resultado: o embedding desses tokens quase nunca recebeu gradiente e permaneceu
próximo da inicialização aleatória. Quando o token aparece num prompt, o modelo
recebe um vetor essencialmente aleatório e responde de forma errática (palavras
sem relação, insultos, evasivas). É o caso extremo dos tokens raros subtreinados:
a lição é que vocabulário e corpus de treinamento precisam estar \textbf{alinhados}.
Há métodos para detectar esses tokens automaticamente (Land \& Bartolo, 2024).
% ---

% --- Solução Q16 ---
\subsection*{XVI — Resposta: (a)}
Cada matriz tem $|V| \times d = 50.000 \times 1.024 = 51.200.000 \approx 51{,}2$M
de parâmetros. Sem weight tying, o modelo mantém duas matrizes independentes desse
tamanho, a tabela de embedding (entrada) e a projeção final \texttt{lm\_head}
(saída), somando $2 \times 51{,}2\text{M} = 102{,}4$M. Com \textbf{weight tying},
a mesma matriz é reutilizada nas duas pontas, eliminando uma das cópias: o total
cai para 51,2M. O GPT-2 usa essa técnica. O custo é relevante sobretudo em modelos
pequenos: no GPT-2 Small ($|V| = 50.257$, $d = 768$), a tabela de embedding sozinha
responde por cerca de um terço dos 117M de parâmetros do modelo.
% ---

\end{document}
